\begin{mistake}{2026-04-12}{01}

\begin{problemcontent}
设三阶矩阵 \(A\)，\(|A|=3\)。将 \(A\) 的第 1、2 行互换，再将第 1 行的 2 倍加到第 2 行，得到矩阵 \(B\)。求 \(|A+B|\)。

\end{problemcontent}

\begin{solutioncontent}
将题述行变换写成左乘初等矩阵。

设交换第 1、2 行对应的初等矩阵为
\[
E_1=\begin{pmatrix}
0&1&0\\
1&0&0\\
0&0&1
\end{pmatrix},
\]
将第 1 行的 2 倍加到第 2 行对应的初等矩阵为
\[
E_2=\begin{pmatrix}
1&0&0\\
2&1&0\\
0&0&1
\end{pmatrix}.
\]
故
\[
B=E_2E_1A=PA,
\qquad
P=E_2E_1=
\begin{pmatrix}
0&1&0\\
1&2&0\\
0&0&1
\end{pmatrix}.
\]

于是
\[
A+B=A+PA=(E+P)A.
\]

从而
\[
|A+B|=|E+P|\,|A|.
\]

计算
\[
E+P=
\begin{pmatrix}
1&1&0\\
1&3&0\\
0&0&2
\end{pmatrix},
\]
故
\[
|E+P|=2\begin{vmatrix}1&1\\1&3\end{vmatrix}=2(3-1)=4.
\]

所以
\[
|A+B|=4\cdot 3=12.
\]

\end{solutioncontent}

\mistakeknowledge{
\begin{itemize}
  \item \textbf{核心公式/定理}：行变换对应左乘初等矩阵；行列式满足 \(|XY|=|X||Y|\)。本题先化为 \(A+B=(E+P)A\) 再求行列式。
  \item \textbf{易错点}：\(|A+B|\neq |A|+|B|\)。行列式对矩阵加法不满足分配律，必须先转化为乘积形式。
  \item \textbf{关联知识}：对矩阵 \(A\) 做行变换时，新矩阵恒可写成 \(PA\)；若做列变换，则应写成 \(AQ\)。左右乘位置不能混淆。
\end{itemize}}

\end{mistake}
